\documentclass[11pt]{article}

\usepackage[margin=1in]{geometry}
\usepackage{amsmath, amssymb, amsthm}
\usepackage{mathtools}
\usepackage{bm}

\DeclareMathOperator*{\clip}{clip}
\DeclareMathOperator{\sign}{sgn}
\newcommand{\Ehat}{\widehat{\mathbb{E}}}
\newcommand{\smean}{s^{\text{MEAN}}}
\newcommand{\ssum}{s^{\text{SUM}}}

\title{Recovery Analysis: \\
Which Normalisation of $s_\tau^\pm,\,s_g^\pm,\,s_{dom}^\pm$ Recovers the Default GRPO Objective at $\delta\to\infty$?}
\author{}
\date{}

\begin{document}
\maketitle

\begin{abstract}
We derive, under verl's \texttt{seq-mean-token-mean} aggregation, the
unclipped limit of BSPO Variants 1 and 2 (simplified and full) for three
candidate definitions of the ratio statistics: (A1) no division anywhere,
(A2) division by $|\tau|$ only at the trajectory level, (A3) mean at every
hierarchy level (our earlier proposal). The recovery criterion is that the
unclipped BSPO objective equals the standard GRPO objective aggregated in
the same way by verl's infra. We find: \textbf{A1 fails}, \textbf{A2 and A3
both recover, for the same reason --- the $/|\tau|$ at the trajectory
level}. Recovery alone therefore does not discriminate A2 from A3; the
choice between them must be made on a separate criterion (practical
clipping sensitivity, hyperparameter count, or numerical behaviour of the
cross-hierarchy subtraction).
\end{abstract}

% =============================================================================
\section{The target: what ``default RL objective'' means under verl's infra}
\label{sec:target}

Per principle~1 (do not change verl's infra), our BSPO aggregation is fixed
to \texttt{seq-mean-token-mean}. Under this mode, the aggregator produces
\[
\mathcal L_{\text{SMTM}}(f) \;=\; \frac{1}{B_{\text{global}}}
\sum_{i=1}^{B}\,\frac{1}{|\tau_i|}\sum_{t=1}^{|\tau_i|}\,f(i,t),
\]
where $f(i,t)$ is the per-token loss tensor (shape $(B,T)$).

The standard GRPO per-token loss is $f_{\text{GRPO}}(i,t)=r_{i,t}\,\hat
A_{\tau_i}$ (advantage broadcast; $r_{i,t}=\pi_\theta(a_t\mid
s_t)/\pi_{\theta_{\text{old}}}(a_t\mid s_t)$). Substituting,
\begin{align*}
\mathcal L_{\text{SMTM}}(r\,\hat A)
  &= \frac{1}{B}\sum_i \hat A_{\tau_i}\,\frac{1}{|\tau_i|}\sum_t r_{i,t} \\
  &= \frac{1}{B}\sum_i \hat A_{\tau_i}\bigl(1 + \smean_{\tau_i}\bigr)
   \quad\text{where }\smean_\tau := \frac{1}{|\tau|}\sum_t (r_t-1) \\
  &= \underbrace{\frac{1}{B}\sum_i \hat A_{\tau_i}}_{=\,0\text{ under GRPO}}
   + \frac{1}{B}\sum_i \hat A_{\tau_i}\,\smean_{\tau_i}.
\end{align*}

GRPO normalisation forces $\sum_{\tau\in g}\hat A_\tau = 0$ for every group
with at least two members; with balanced groups across $B$, the first sum
is exactly zero. Define
\[
\boxed{\quad J_\star \;:=\; \Ehat_\tau\bigl[\smean_\tau \cdot \hat A_\tau\bigr]
        \;=\; \tfrac{1}{B}\sum_i \smean_{\tau_i}\,\hat A_{\tau_i}. \quad}
\]
This is the target: the unclipped BSPO objective must equal $J_\star$ at
$\delta\to\infty$.

The key observation is that $J_\star$ involves $\smean_\tau$, the
\emph{trajectory mean} of the per-token ratio deviation. The $/|\tau|$ is
built into the target by verl's aggregator itself.

% =============================================================================
\section{The three candidate normalisations}

All three share the token-level deviation construction; they differ only
in where the normalisers are placed.

\paragraph{A1 (no division anywhere --- ``remove $/|\tau|$'').}
\begin{align*}
s_\tau^{\pm,\text{A1}}    &= \sum_{t\in\tau}\max\bigl(\pm(r_t-1),\,0\bigr) \\
s_g^{\pm,\text{A1}}       &= \sum_{\tau\in g}\max\bigl(\pm s_\tau^{\text{A1}},\,0\bigr) \\
s_{dom}^{\pm,\text{A1}}   &= \sum_{g\in dom}\max\bigl(\pm s_g^{\text{A1}},\,0\bigr)
\end{align*}

\paragraph{A2 (\textit{the only} normaliser is $/|\tau|$ at the trajectory level).}
\begin{align*}
s_\tau^{\pm,\text{A2}}    &= \frac{1}{|\tau|}\sum_{t\in\tau}\max\bigl(\pm(r_t-1),\,0\bigr) \\
s_g^{\pm,\text{A2}}       &= \sum_{\tau\in g}\max\bigl(\pm s_\tau^{\text{A2}},\,0\bigr) \\
s_{dom}^{\pm,\text{A2}}   &= \sum_{g\in dom}\max\bigl(\pm s_g^{\text{A2}},\,0\bigr)
\end{align*}

\paragraph{A3 (mean at every level --- the section~3 edit in \texttt{bisim\_eqs\_multidomain.tex}).}
\begin{align*}
s_\tau^{\pm,\text{A3}}    &= \frac{1}{|\tau|}\sum_{t\in\tau}\max\bigl(\pm(r_t-1),\,0\bigr) \\
s_g^{\pm,\text{A3}}       &= \frac{1}{|g|}\sum_{\tau\in g}\max\bigl(\pm s_\tau^{\text{A3}},\,0\bigr) \\
s_{dom}^{\pm,\text{A3}}   &= \frac{1}{|dom|}\sum_{g\in dom}\max\bigl(\pm s_g^{\text{A3}},\,0\bigr)
\end{align*}

In all three, $s_z = s_z^+ - s_z^-$ by the identity $\max(x,0)-\max(-x,0)=x$.

% =============================================================================
\section{Useful identities}
\label{sec:identities}

The following are purely algebraic; they hold for all three variants.

\begin{lemma}[Trajectory identity]
\label{lem:tau}
For any definition of $s_\tau^\pm$ that uses $\max(\pm x,\,0)$ decomposition,
$s_\tau \;=\; \sum_t\alpha_\tau\,(r_t-1)$, with the prefactor
$\alpha_\tau = 1$ under A1 and $\alpha_\tau = 1/|\tau|$ under A2/A3.
In particular, $s_\tau^{\text{A1}} = |\tau|\,\smean_\tau$ and
$s_\tau^{\text{A2}} = s_\tau^{\text{A3}} = \smean_\tau$.
\end{lemma}

\begin{lemma}[Group-/domain-level $s_z$ are constants on their support]
\label{lem:const}
$s_g$ depends only on $g$ (not on $\tau$ within $g$). $s_{dom}$ depends
only on $dom$. This holds under all three normalisations.
\end{lemma}

\begin{lemma}[GRPO zero-group-mean]
\label{lem:grpo}
Under GRPO trajectory-level advantage normalisation with group size
$\geq 2$,
\[
\sum_{\tau\in g}\hat A_\tau \;=\; 0 \quad\text{for every }g.
\]
Consequently, if $\hat A_g := (1/|g|)\sum_{\tau\in g}\hat A_\tau$ and
$\hat A_{dom}$ is any linear aggregation of $\{\hat A_g\}$, then
$\hat A_g = 0$ and $\hat A_{dom}=0$.
\end{lemma}

\begin{lemma}[Higher-level constants vanish against $\hat A_\tau$]
\label{lem:vanish}
For any group-level constant $c_g$ and any domain-level constant $c_{dom}$,
\[
\Ehat_\tau\bigl[c_{g(\tau)}\,\hat A_\tau\bigr] \;=\; 0,
\qquad
\Ehat_\tau\bigl[c_{dom(\tau)}\,\hat A_\tau\bigr] \;=\; 0,
\]
where $g(\tau)$ is the group containing $\tau$ and $dom(\tau)$ is its domain.
\end{lemma}
\begin{proof}
$\Ehat_\tau[c_g\,\hat A_\tau] = \tfrac{1}{B}\sum_g\sum_{\tau\in g}
c_g\,\hat A_\tau = \tfrac{1}{B}\sum_g c_g\sum_{\tau\in g}\hat A_\tau = 0$
by Lemma~\ref{lem:grpo}. Similarly for $c_{dom}$.
\end{proof}

% =============================================================================
\section{Variant 1 unclipped}

Variant 1 simplest: $\hat J_1 = \Ehat_\tau\bigl[\Delta_\tau^{(\text{reg})}\,\hat A_\tau\bigr]$.
At $\delta\to\infty$, $\clip(s_\tau^{\pm},\delta)\to s_\tau^{\pm}$ and so
$\Delta_\tau^{(\text{reg})} \to s_\tau^+ - s_\tau^- = s_\tau$.

Thus $\hat J_1^{(\text{unclipped})} = \Ehat_\tau[s_\tau\,\hat A_\tau]$.

Substituting Lemma~\ref{lem:tau}:

\begin{center}
\begin{tabular}{l l l}
\textbf{Approach} & \textbf{$s_\tau$} & \textbf{$\hat J_1^{(\infty)}$} \\ \hline
A1  & $|\tau|\,\smean_\tau$ & $\Ehat_\tau\bigl[|\tau|\,\smean_\tau\,\hat A_\tau\bigr]$ \\
A2  & $\smean_\tau$ & $\Ehat_\tau\bigl[\smean_\tau\,\hat A_\tau\bigr] \;=\; J_\star$ \\
A3  & $\smean_\tau$ & $\Ehat_\tau\bigl[\smean_\tau\,\hat A_\tau\bigr] \;=\; J_\star$ \\
\end{tabular}
\end{center}

A1 produces the wrong $\smean_\tau\cdot\hat A_\tau$ \emph{weighted by
$|\tau|$}. For $|\tau|$ uniform across the batch this is a constant
rescaling (benign); for mixed-length batches (always the case in practice)
this is not $J_\star$. \textbf{A1 fails Variant~1 recovery.}

A2 and A3 recover $J_\star$ identically because they place the same
$/|\tau|$ at the trajectory level.

% =============================================================================
\section{Variant 2 unclipped (simplified form)}

Simplified Variant 2: $\hat J_2 = \Ehat_\tau\bigl[(\Delta_\tau^{(\text{reg})} - \Delta_g^{(\text{reg})})\,\hat A_\tau\bigr]$.
At $\delta\to\infty$, $\Delta_\tau^{(\text{reg})}\to s_\tau$ and
$\Delta_g^{(\text{reg})}\to s_g$ (both with the signs of $s$).

\[
\hat J_2^{(\text{unclipped})}
  = \Ehat_\tau\bigl[s_\tau\,\hat A_\tau\bigr]
   - \Ehat_\tau\bigl[s_{g(\tau)}\,\hat A_\tau\bigr].
\]

By Lemma~\ref{lem:const} $s_g$ is a group-level constant; by
Lemma~\ref{lem:vanish} its product with $\hat A_\tau$ averages to zero
under GRPO \emph{regardless of how $s_g$ is normalised}.

\[
\hat J_2^{(\text{unclipped})} \;=\; \Ehat_\tau[s_\tau\,\hat A_\tau] \;-\; 0.
\]

Recovery status is therefore identical to Variant~1:

\begin{center}
\begin{tabular}{l l}
\textbf{Approach} & \textbf{$\hat J_2^{(\infty)}$} \\ \hline
A1 & $\Ehat_\tau\bigl[|\tau|\,\smean_\tau\,\hat A_\tau\bigr]$ (fails) \\
A2 & $J_\star$ (recovers) \\
A3 & $J_\star$ (recovers) \\
\end{tabular}
\end{center}

\textbf{Crucial observation.} The normalisation choice on $s_g$ (sum vs.\
mean over $\tau\in g$) has \emph{no effect} on the unclipped recovery. Both
A2 and A3 succeed, and they succeed for exactly the same reason: $s_g$ is
a group-level constant, and any group-level constant pushed through
$\Ehat_\tau$ with $\hat A_\tau$ vanishes under GRPO's zero-group-mean
property. The absolute scale of $s_g$ is irrelevant to this identity.

% =============================================================================
\section{Variant 2 unclipped (full hierarchical form)}

Full Variant 2 before telescoping:
\[
\hat J_2^{\text{full}}
  = \Ehat_\tau\bigl[\Delta_{dom}^{(\text{reg})}\,\hat A_{dom}
       + (\Delta_g^{(\text{reg})}-\Delta_{dom}^{(\text{reg})})\,\hat A_g
       + (\Delta_\tau^{(\text{reg})}-\Delta_g^{(\text{reg})})\,\hat A_\tau\bigr].
\]

At $\delta\to\infty$ and using Lemma~\ref{lem:grpo} ($\hat A_g = \hat A_{dom} = 0$):
\begin{align*}
\hat J_2^{\text{full}\,(\text{unclipped})}
  &= \Ehat_\tau\bigl[s_{dom}\cdot 0 + (s_g-s_{dom})\cdot 0 + (s_\tau-s_g)\,\hat A_\tau\bigr] \\
  &= \Ehat_\tau\bigl[(s_\tau-s_g)\,\hat A_\tau\bigr] \\
  &= \Ehat_\tau[s_\tau\,\hat A_\tau] - 0
   \qquad\text{(Lemma~\ref{lem:vanish})}.
\end{align*}

Recovery status identical to the simplified form. A1 fails; A2 and A3
recover.

% =============================================================================
\section{Why the normalisation at $g$ and $dom$ is invisible to the recovery criterion}

The identities in \S\ref{sec:identities} make it explicit:

\begin{enumerate}
\item $s_g$ and $s_{dom}$ appear in the unclipped objective \emph{only
multiplied against $\hat A_\tau$, $\hat A_g$, or $\hat A_{dom}$}.
\item All three advantage terms vanish in expectation under GRPO
(Lemmas~\ref{lem:grpo}, \ref{lem:vanish}).
\item The magnitude of $s_g$ and $s_{dom}$ is therefore absorbed into a
multiplication-by-zero at the expectation level.
\end{enumerate}

As a result, \textbf{the unclipped-recovery criterion cannot distinguish
between A2 and A3}. It only distinguishes them from A1 (which fails because
$s_\tau$ itself acquires a spurious $|\tau|$ factor that does \emph{not}
vanish against $\hat A_\tau$).

This answers the user's original concern: ``can we still recover the
default GRPO objective when $s_g$ and $s_{dom}$ are at different scales?''
\textbf{Yes, for the same structural reason in both A2 and A3.} The
trajectory-level $/|\tau|$ is the load-bearing ingredient; the higher-level
normalisations only matter for finite-$\delta$ clipping behaviour.

% =============================================================================
\section{The canonical division factor determined by the recovery rule}

Collecting the derivations: among the three candidates considered, the
\emph{minimum requirement} imposed by the recovery rule is

\[
\boxed{\quad
s_\tau^\pm \text{ must contain a factor of }\tfrac{1}{|\tau|}\text{ inside the one-sided token sums.}
\quad}
\]

Higher-level normalisations are \emph{undetermined} by recovery alone.
Formally, the recovery rule admits the entire family
\[
s_z^{\pm,\alpha_z}
  \;=\; \alpha_z \sum_{\text{members of }z}\max\bigl(\pm s_{\text{child}},\,0\bigr),
  \qquad z\in\{g, dom\},
\]
for any positive weights $\{\alpha_g,\alpha_{dom}\}$. A2 picks
$\alpha_g=\alpha_{dom}=1$ (pure sums); A3 picks
$\alpha_g=1/|g|,\alpha_{dom}=1/|dom|$ (means). Both satisfy recovery.

% =============================================================================
\section{Separating criteria (not recovery): what differs between A2 and A3}

Since recovery does not discriminate, we summarise the practical
differences for a downstream choice:

\begin{center}
\renewcommand{\arraystretch}{1.2}
\begin{tabular}{l l l}
\textbf{Property} & \textbf{A2} & \textbf{A3} \\ \hline
$s_\tau$ scale        & $O(\varepsilon)$                & $O(\varepsilon)$ \\
$s_g$ scale           & $O(|g|\cdot\varepsilon)$        & $O(\varepsilon)$ \\
$s_{dom}$ scale       & $O(|dom|\cdot|g|\cdot\varepsilon)$ & $O(\varepsilon)$ \\
$\delta$ hyperparam count (if all levels must not saturate) & 3 distinct scales & 1 scale \\
Reduction primitive in verl       & custom \texttt{scatter\_sum}   & \texttt{scatter\_mean} (already used) \\
Value of $\Delta_\tau^{\text{(reg)}}-\Delta_g^{\text{(reg)}}$ at finite $\delta$ & dominated by $\Delta_g$ when $\delta_g < |g|\,\varepsilon$ & deviation-from-mean, bounded by $\delta$ \\
\end{tabular}
\end{center}

The user's own characterisation (``token ratio naturally has much higher
average magnitude than trajectory agg ratio \ldots\ we anticipate to
subtract the smaller scale from higher part of the hierarchy'') matches
the A2 scale cascade only in the direction \emph{higher-hierarchy is
larger, not smaller}. Under A2, $|s_g|\geq|s_\tau|$; under A3, all three
are comparable. If the stated preference is ``higher hierarchy holds a
smaller-scale correction that is \emph{subtracted from} the lower-level
quantity,'' this is realised exactly by A3 under the finite-batch
variance interpretation: $s_g$ is a $1/\sqrt{|g|}$-contracted estimator
of the group mean of $s_\tau$, and $s_\tau - s_g$ is a mean-centred
deviation whose magnitude is bounded by the within-group spread of
$s_\tau$ rather than by $|s_\tau|$ itself.

% =============================================================================
\section{Explicit failure of A1: worked micro-example}

Take a batch of two trajectories:
\begin{itemize}
\item $\tau_1$: $|\tau_1|=3$, ratios $[1.1,\;1.0,\;1.2]$, so $\smean_{\tau_1}=0.1$.
\item $\tau_2$: $|\tau_2|=300$, ratios all $1.001$ except one token at
      $1.0$, so $\smean_{\tau_2}\approx 0.001$.
\end{itemize}
Assume for simplicity one group containing both, with GRPO advantages
$\hat A_{\tau_1}=+1,\ \hat A_{\tau_2}=-1$ (satisfying zero-group-mean).

\noindent\textbf{Target.}
$J_\star = \tfrac{1}{2}\bigl(0.1\cdot 1 + 0.001\cdot(-1)\bigr) = 0.0495$.

\noindent\textbf{A1} ($s_\tau^{\text{A1}}=|\tau|\smean$): \;
$s_{\tau_1}=0.3,\ s_{\tau_2}\approx 0.3$. \\
$\hat J_1^{(\text{A1, unclipped})}
= \tfrac{1}{2}(0.3\cdot 1 + 0.3\cdot(-1)) = 0$.

\noindent\textbf{A2 / A3} ($s_\tau=\smean_\tau$): \;
$\hat J_1^{(\text{A2/A3, unclipped})}
= \tfrac{1}{2}(0.1\cdot 1 + 0.001\cdot(-1)) = 0.0495 = J_\star$.

A1 returns $0$ for an example where $J_\star=0.0495$: a $50\times$ error
in relative terms, purely from the $|\tau|$-weighting bias. A2 and A3
return the exact target.

% =============================================================================
\section{Conclusion}

\begin{itemize}
\item \textbf{A1 is ruled out by recovery.} The $/|\tau|$ at the
      trajectory level is load-bearing and cannot be removed.
\item \textbf{A2 and A3 both satisfy recovery.} The higher-level
      normalisation (sum vs.\ mean) is invisible to the unclipped-limit
      criterion.
\item \textbf{The canonical factor determined by recovery alone is
      $1/|\tau|$ on $s_\tau^\pm$, nothing else.}
\item \textbf{To choose between A2 and A3 one must invoke a criterion
      beyond recovery}, e.g.\ clipping sensitivity at finite $\delta$,
      hyperparameter count, subtraction magnitude, or alignment with
      verl's existing \texttt{scatter\_mean} primitive. On those grounds
      A3 is preferable, but the recovery property itself does not force
      it.
\end{itemize}

\end{document}
